\documentclass[10pt,a4paper]{amsart}
\usepackage[margin=19mm]{geometry}
\usepackage{amsmath,amssymb,mathtools}
\usepackage[scr=rsfs,cal=euler]{mathalfa}
\usepackage{xfrac}
\usepackage[unicode,hidelinks,bookmarks=false]{hyperref}
\newcommand{\ie}{i.e.\ }
\newcommand{\cf}{cf.\ }
\newcommand{\resp}{respectively\ }
\newcommand{\Subsection}{\subsection}
\providecommand{\nobreakdash}{\nobreak\hbox{-}}
\setlength{\emergencystretch}{2em}
\multlinegap0pt
\usepackage{amssymb}
\usepackage{mathtools}
\usepackage[shortlabels]{enumitem}
\usepackage{xcolor}
\newtheorem{theorem}{Theorem}
\newcommand{\dd}{\,\mathrm{d}}
\title{Positive rational series for reciprocal powers of Catalan's constant and Dirichlet beta values}
\author{Narendra Bhandari}
\date{Source: arXiv:2608.00429v1 (CC BY 4.0)}
\begin{document}
\maketitle

\noindent\emph{Source and licence.} Positive rational series for reciprocal powers of Catalan's constant and Dirichlet beta values. Version arXiv:2608.00429v1 (CC BY 4.0), distributed by arXiv under the Creative Commons Attribution 4.0 International licence, http://creativecommons.org/licenses/by/4.0/, as recorded in the arXiv licence metadata for this exact version and shown on the abstract page https://arxiv.org/abs/2608.00429v1. Full licence notice: \texttt{licenses/arxiv-2608.00429-CC-BY-4.0.txt}.\par\smallskip

\noindent\textit{Editorial scope.} Counted unit: the article's theorem thm:H-coefficients (source0.tex lines 435--458). The article's complete original proof of it is delivered with it (source0.tex lines 460--559, one \texttt{proof} environment of 100 lines). No mathematics is written, repaired or re-derived: the statement and the proof are the article's own lines, in the article's own notation, and no retained line is substituted or re-flowed.  Retained uncounted as the source-local context the proof needs: lines 311--317; lines 326--331; lines 372--374; lines 432--434.  Nothing was omitted from any retained span, so the package carries no omission record.  No retained line contains a citation, so the wrapper carries no bibliography and no bibliography entry.  No retained line contains a figure environment, an image-inclusion command or drawing code, so no image is delivered and the package declares no asset dependency.  Editorial formatting only, in addition: an AMS \texttt{amsart} preamble that restates the article's own declarations and macros used by the retained lines, namely the environments \texttt{theorem} on separate counters, together with the standard packages those lines need.  The retained environments print in this document as Theorem 1 in order of appearance, whereas the article prints the same items with the numbering of its own full document.  The complete native source of the article as distributed in the arXiv e-print for this exact version is bundled unchanged as \texttt{sources/206539/source0.tex}.
\begin{equation}\label{eq:F-integral}
 F_s(t)
 =
 \frac1{\Gamma(s)}
 \int_0^1
 \frac{(-\log x)^{s-1}}{1+tx^2}\dd x.
\end{equation}

\begin{equation}\label{eq:moment}
 \frac1{(2k+1)^s}
 =
 \frac1{\Gamma(s)}
 \int_0^1x^{2k}(-\log x)^{s-1}\dd x.
\end{equation}

\begin{align}\label{eq-gamma}
    \int_0^1(-\log x)^{s-1}\dd x=\Gamma(s)<\infty.
\end{align}

\begin{equation}\label{eq:H-definition}
 H_s(w)=F_s\!\left(\frac{w}{1-w}\right).
\end{equation}

\begin{theorem}\label{thm:H-coefficients}
For real \(s>0\) and \(|w|<1\),
\begin{equation}\label{eq:H-expansion}
 H_s(w)=1-\sum_{n=1}^{\infty}q_{s,n}w^n,
\end{equation}
where
\begin{align}
 q_{s,n}
 &=
 \sum_{k=1}^{n}
 \frac{(-1)^{k+1}}{(2k+1)^s}
 \binom{n-1}{k-1},                       \label{eq:q-finite}\\
 &=
 \frac1{\Gamma(s)}
 \int_0^1
 x^2(1-x^2)^{n-1}
 (-\log x)^{s-1}\dd x.                 \label{eq:q-integral}
\end{align}
In particular,
\(
 q_{s,n}>0
 \) for \(s>0,\ n\geq1.
\)
\end{theorem}

\begin{proof}
Substituting \eqref{eq:H-definition} into the integral representation
\eqref{eq:F-integral} gives
\[
 \begin{aligned}
 H_s(w)
 &=
 \frac1{\Gamma(s)}
 \int_0^1
 \frac{(1-w)(-\log x)^{s-1}}
 {1-w+x^2w}\dd x\\
 &=
 \frac1{\Gamma(s)}
 \int_0^1
 \frac{(1-w)(-\log x)^{s-1}}
 {1-(1-x^2)w}\dd x.
 \end{aligned}
\]

For \(|w|<1\), expanding the denominator as a geometric series gives
\[
 \begin{aligned}
 \frac{1-w}{1-(1-x^2)w}
 &=
 1-\frac{x^2w}{1-(1-x^2)w}=
 1-x^2w
 \sum_{m=0}^{\infty}(1-x^2)^mw^m\\
 &=
 1-\sum_{n=1}^{\infty}
 x^2(1-x^2)^{n-1}w^n.
 \end{aligned}
\] 
To justify integration term by term, fix \(|w|\leq\rho<1\).  Then
\[
 \begin{aligned}
 \sum_{n=1}^{\infty}
 \left|
 x^2(1-x^2)^{n-1}w^n
 \right|
 &\leq
 \sum_{n=1}^{\infty}\rho^n=
 \frac{\rho}{1-\rho}.
 \end{aligned}
\]
Together with \eqref{eq-gamma}, this bound justifies integration term by
term for \(|w|\leq\rho<1\).  We therefore obtain

\[
 \begin{aligned}
 H_s(w)
 &=
 \frac1{\Gamma(s)}
 \int_0^1(-\log x)^{s-1}\dd x-
 \sum_{n=1}^{\infty}
 \left[
 \frac1{\Gamma(s)}
 \int_0^1
 x^2(1-x^2)^{n-1}
 (-\log x)^{s-1}\dd x
 \right]w^n.
 \end{aligned}
\]
Comparison with \eqref{eq:H-expansion}, using
\eqref{eq-gamma}, now yields
\[
 q_{s,n}
 =
 \frac1{\Gamma(s)}
 \int_0^1
 x^2(1-x^2)^{n-1}
 (-\log x)^{s-1}\dd x,
\]
This is \eqref{eq:q-integral}.  To recover the finite formula, expand
\[
 (1-x^2)^{n-1}
 =
 \sum_{j=0}^{n-1}
 (-1)^j\binom{n-1}{j}x^{2j}.
\]
Substituting this finite expansion into \eqref{eq:q-integral} and
using \eqref{eq:moment} gives
\[
 \begin{aligned}
 q_{s,n}
 &=
 \frac1{\Gamma(s)}
 \sum_{j=0}^{n-1}
 (-1)^j\binom{n-1}{j}
 \int_0^1
 x^{2j+2}(-\log x)^{s-1}\dd x\\
 &=
 \sum_{j=0}^{n-1}
 \frac{(-1)^j}{(2j+3)^s}
 \binom{n-1}{j}.
 \end{aligned}
\]
Relabeling \(k=j+1\) gives \eqref{eq:q-finite}.  Positivity is
immediate from \eqref{eq:q-integral}: every factor in the integrand is
positive on \(0<x<1\), and \(\Gamma(s)>0\).
\end{proof}

\end{document}
